\definecolor{nodefill}{HTML}{DFE3E8}%
\tikzset{%
  rnode/.style={circle, draw, thick, minimum size=5mm, inner sep=1pt,
                font=\scriptsize\bfseries, fill=nodefill},%
  rlbl/.style={font=\tiny, midway, above},%
  rlblb/.style={font=\tiny, midway, below},%
  missing/.style={dashed, gray},%
  missingnode/.style={rnode, dashed, gray, text=gray, fill=none},%
}%

\begin{tabularx}{\textwidth}{p{2.2cm}cX}
\toprule
Group & Graph & Description \\
\midrule
\multirow{4}{=}{Hypothesis handling} & \begin{tikzpicture}[>=Stealth, baseline=(current bounding box.center), node distance=7mm]
  \node[rnode] (e) {E};
  \node[rnode, right=of e] (j) {J};
  \node[right=4mm of j, draw=none, font=\scriptsize] (dots) {\ldots};
  \node[rnode, right=4mm of dots] (h) {H};
  \node[rnode, right=of h] (t) {T};
  \draw[->] (e) -- node[rlbl] {inf.} (j);
  \draw[->] (h) -- node[rlbl] {tests} (t);
\end{tikzpicture} & \textbf{Evidence led hypothesis generation}. Evidence is observed first; a hypothesis is formed afterward \\
 & \begin{tikzpicture}[>=Stealth, baseline=(current bounding box.center), node distance=7mm]
  \node[rnode] (h1) {H$_1$};
  \node[rnode, right=15mm of h1] (h2) {H$_2$};
  \node[rnode, below left=5mm and 0mm of h1] (t1) {T};
  \node[rnode, below right=5mm and 0mm of h2] (t2) {T};
  \draw[<->, dashed] (h1) -- node[rlbl] {competes} (h2);
  \draw[->] (h1) -- node[left, font=\tiny] {tests} (t1);
  \draw[->] (h2) -- node[right, font=\tiny] {tests} (t2);
\end{tikzpicture} & \textbf{Hypothesis reranking}. Competing hypotheses are compared as new evidence arrives \\
 & \begin{tikzpicture}[>=Stealth, baseline=(current bounding box.center), node distance=7mm]
  \node[rnode] (h1) {H};
  \node[rnode, right=of h1] (t) {T};
  \node[rnode, right=of t] (e) {E};
  \node[rnode, right=of e] (j) {J};
  \node[rnode, right=of j] (h2) {H$_2$};
  \draw[->] (h1) -- node[rlbl] {tests} (t);
  \draw[->] (t) -- node[rlbl] {obs.} (e);
  \draw[->] (e) -- node[rlbl] {inf.} (j);
  \draw[->] (h1) to[bend right=40] node[rlblb] {upd.} (h2);
\end{tikzpicture} & \textbf{Refutation driven belief revision}. Evidence triggers a belief update to a new hypothesis \\
 & \begin{tikzpicture}[>=Stealth, baseline=(current bounding box.center), node distance=7mm]
  \node[rnode] (t1) {T};
  \node[rnode, right=of t1] (e) {E};
  \node[right=4mm of e, draw=none, font=\scriptsize] (dots) {\ldots};
  \node[rnode, right=4mm of dots] (h) {H};
  \node[rnode, right=of h] (t2) {T};
  \draw[->] (t1) -- node[rlbl] {obs.} (e);
  \draw[->] (h) -- node[rlbl] {tests} (t2);
\end{tikzpicture} & \textbf{Explore then test transition}. Exploration precedes hypothesis formation, which then drives testing \\
\midrule[0.015em]
\multirow{1}{=}{Evidence handling} & \begin{tikzpicture}[>=Stealth, baseline=(current bounding box.center), node distance=7mm]
  \node[rnode] (h) {H};
  \node[rnode, right=10mm of h, yshift=7mm] (t1) {T$_1$};
  \node[rnode, right=10mm of h] (t2) {T$_2$};
  \node[rnode, right=10mm of h, yshift=-7mm] (t3) {T$_3$};
  \node[rnode, right=10mm of t2] (e) {E};
  \draw[->] (h) -- (t1);
  \draw[->] (h) -- (t2);
  \draw[->] (h) -- (t3);
  \draw[->] (t1) -- (e);
  \draw[->] (t2) -- (e);
  \draw[->] (t3) -- (e);
\end{tikzpicture} & \textbf{Convergent multi test evidence}. One hypothesis is evaluated via multiple independent tests \\
\midrule[0.015em]
\multirow{2}{=}{Inquiry control} & \begin{tikzpicture}[>=Stealth, baseline=(current bounding box.center), node distance=7mm]
  \node[rnode] (h) {H};
  \node[rnode, right=of h] (t) {T};
  \node[rnode, right=of t] (e) {E};
  \node[rnode, right=of e] (j) {J};
  \draw[->] (h) -- node[rlbl] {tests} (t);
  \draw[->] (t) -- node[rlbl] {obs.} (e);
  \draw[->] (e) -- node[rlbl] {inf.} (j);
  \draw[->, bend left=50] (j) to node[rlblb] {tests} (t);
\end{tikzpicture} & \textbf{Fixed hypothesis test tuning}. Hypothesis is held fixed while tests are iteratively adjusted \\
 & \begin{tikzpicture}[>=Stealth, baseline=(current bounding box.center), node distance=7mm]
  \node[rnode] (j) {J};
  \node[rnode, right=of j] (t) {T};
  \node[rnode, right=of t] (e) {E};
  \draw[->] (j) -- node[rlbl] {tests} (t);
  \draw[->] (t) -- node[rlbl] {obs.} (e);
\end{tikzpicture} & \textbf{Evidence guided test redesign}. A judgment motivates a new test, which then produces new evidence \\
\bottomrule
\end{tabularx}